\section{The new finite degrees in the cyclotomic tower}
\label{sec:component-extension}

\begin{corollary}[KT-COMPONENT-EXTENSION]
Use the notation of K17-COMPONENT-CONTROL. For every odd prime
$p$ and positive integer $k$ whose residue modulo $(p-1)/2$
has a representative $j$ with
\[
                         -16\le j\le15,
\]
one has
\[
 K_{4k}(\mathbf Z)\{p\}=0,\qquad
 A_n^{[-2k]}=0\ (n\ge0),\qquad X^{[-2k]}=0.
\]
For $p\ge5$, the actual ramified component $M^{[2k+1]}$ is
integrally free of rank one over the Iwasawa algebra $\Lambda$.
For every odd $p$, the full-domain Soul\'e characters in weights
$27$, $29$ and $31$ are onto in the precise cyclotomic-product
normalization of K17, (K19).
\end{corollary}
\begin{proof}
The carried K17-TWELVE-DEGREES-PERIOD supplies representatives
$-16\le j\le12$. Theorem KT-K52-K56-K60-VANISH supplies $j=13,14,15$.
K17's exact comparison (K2) gives $A_0^{[-2j]}=0$ in these
additional degrees. Congruent exponents select the same
$\mathbf F_p[\operatorname{Gal}(\mathbf Q(\mu_p)/\mathbf Q)]$
eigenspace, so (K2) gives the primary vanishing for the stated
residue classes. K17-COMPONENT-CONTROL, beginning at line~10975
of the supplied paper, then gives every finite-layer and
inverse-limit vanishing and the integral rank-one freeness.
For the exceptional odd characters $\omega^{\pm1}$, its cited
K17-ODD-COMPANION-EXCLUSION, beginning at line~11084,
supplies the same conclusion. These are carried control
statements, not a new Iwasawa-control proof.

Finally the full-domain index formula K17 (K18), in the section
beginning before line~8390, gives
\[
 [\mathbf Z_p(m):\kappa_m(N_p)]
       =|K_{2m-2}(\mathbf Z)\{p\}|.
\]
For $m=27,29,31$ the right sides are one. This proves the three
surjectivity claims with exactly its domain and normalization.
No assertion about the image of a deeper filtration subgroup
follows from this formula.
\end{proof}

The result enlarges selected character ranges in LM6 A.44--A.45
and the bridges Q06; it does not exhaust the characters at a
large prime or establish the full Ihara kernel assertion.
The periodic statement is about vanishing, not equality of the
orders or an isomorphism of all $K$-groups in congruent degrees.
